\section{Method}
\label{sec:method}

As shown in Fig.~\ref{fig:main}, our system uses the 3D motion coefficients as the intermediate representation for talking head generation. We first extract the coefficients from the original image. Then, the realistic 3DMM motion coefficients are generated by ExpNet and PoseVAE individually. Finally, a 3D-aware face render is proposed to produce the talking head videos. Below, we give a brief introduction to the 3D face model as preliminaries in Sec.~\ref{sec:preliminary}, the audio-driven motion coefficients generation and the coefficients-driven image animator we design in Sec.~\ref{sec:audio2coeff} and Sec.~\ref{sec:render}, respectively.  

\begin{figure}[tp]
    \centering
    \includegraphics[width=\columnwidth]{figs/images/main_of_sadtalker.pdf}
    \vspace{-2em}
    \caption{Main pipeline. Our method uses the coefficients of 3DMM as intermediate motion representation. To this end, we first generate realistic 3D motion coefficients~(facial expression $\beta$, head pose $\rho$) from audio, then these coefficients are used to implicitly modulate the 3D-aware face render for final video generation.}
    \label{fig:main}
\end{figure}



\subsection{Preliminary of 3D Face Model}
\label{sec:preliminary}
3D information is crucial to improve the realness of the generated video since the real video is captured in the 3D environment. However, previous works~\cite{pcavs, wav2lip, zhou2020makelttalk} have rarely been a consideration in 3D space since it is hard to obtain accurate 3D coefficients from a single image and the high-quality face render is also hard to design. 
Inspired by the recent single image deep 3D reconstruction method~\cite{deng2019accurate}, we consider the space of the predicted 3D Morphable Models~(3DMMs) as our intermediate representation.
In 3DMM, the 3D face shape $\mathbf{S}$ can be decoupled as:
\begin{equation}
\mathbf{S} = \overline{\mathbf{S}} + \mathbf{\alpha} \mathbf{U}_{id} + \mathbf{\beta} \mathbf{U}_{exp},
\end{equation}
where $\overline{\mathbf{S}}$ is the average shape of the 3D face, $\mathbf{U}_{id}$ and $\mathbf{U}_{exp}$ are the orthonormal basis of identity and expression of LSFM morphable model~\cite{3dmm}. 
Coefficients $\mathbf{\alpha} \in \mathbb{R}^{80}$ and $\mathbf{\beta} \in \mathbb{R}^{64}$ describe the person identity and expression, respectively. 
To preserve pose variance, coefficients $\mathbf{r} \in SO(3)$ and $\mathbf{t} \in \mathbb{R}^{3}$ denote the head rotation and translation. To achieve identity irrelevant coefficients generation~\cite{ren2021pirenderer}, we only model the parameters of motion as $ \left \{ \mathbf{\beta}, \mathbf{r}, \mathbf{t} \right \}$. We learn the head pose $ \rho = [\mathbf{r}, \mathbf{t}]$ and expression coefficients $\beta$ individually from the driving audio as introduced before. Then, these motion coefficients are used to implicitly modulate our face render for final video synthesis. 
% A similar approach has also be used in PIRenderer~\cite{ren2021pirenderer}, comparing with their method, we remove the crop coefficient from their method to achieve realistic. we are focusing on generating the realistic 3D coefficients since their method focuses on the face rendering.
% Then, we can model the motion of the driving face with a parameter set $\mathbf{p} = \left \{ \mathbf{\beta}, \mathbf{r}, \mathbf{t} \right \}$ extracted by an existing 3D face reconstruction model~\cite{deng2019accurate}.

\subsection{Motion Coefficients Generation through Audio}
\label{sec:audio2coeff}
As introduced above, the 3D motion coefficients contain both head pose and expression where the head pose is a global motion and the expression is relatively local. To this end, learning everything altogether will cause huge uncertainty in the network since the head pose has a relatively weak relationship with audio while the lip motion is highly connected.
We generate the motion of the head pose and expression using the proposed PoseVAE and ExpNet, respectively introduced below. 

% The coefficients of the source image is a good reference to generate realistic coefficients of the 3D model. For the input face image, we extract the initial 3d head pose $\rho_{0}$ and expression coefficients $\beta_{0}$ from a pre-trained 3D face reconstruction model~\cite{deng2019accurate}. As for the audio, we pre-process the audio to mel-spectrum feature $a_{\{1...n\}}$~($n$ is the total number of the audio window) \xiaodong{add audio pre-processing}. 

\paragraph{ExpNet}

Learning a generic model which produces accurate expression coefficients from audio is extremely hard for two reasons: 1) audio-to-expression is not a one-to-one mapping task for different identities. 2) there are some audio-irrelevant motions in the expression coefficients and it will influence the prediction's accuracy. 
%our method design
Our ExpNet is designed to reduce these uncertainties. As for the identity issue, we connect the expression motion to the specific person via the first frame's expression coefficients $\beta_0$. 
To reduce the motion weight of other facial components in natural talking, we use the \textit{lip motion only} coefficients as the coefficient target through the pre-trained network of Wav2Lip~\cite{wav2lip} and deep 3D reconstruction~\cite{deng2019accurate}. Then, other minor facial motions~(\eg, eye blink) can be leveraged via the additional landmark loss on the rendered images.

\begin{figure}[tp]
    \centering
    \includegraphics[width=\columnwidth]{figs/images/expnet_v2.pdf}
    \vspace{-2.5em}
    \caption{The structure of our ExpNet. We involve a monocular 3D face reconstruction model~\cite{deng2019accurate}~($R_e$ and $R_d$) to learn the realistic expression coefficients. Where $R_e$ is a pretrained 3DMM coefficients estimator and $R_d$ is a differentiable 3D face render without learnable parameters. We use the reference expression $\beta_0$ to reduce the uncertainty of identity and the generated frame from pre-trained Wav2Lip~\cite{wav2lip} and the first frame as target expression coefficients since it only contains the lip-related motions. }
    \label{fig:expnet}
\end{figure}

%condition on audio, the reference expression $\beta_{0}$ from the first frame and the controllable blink factor $z_{blink}$

As shown in Figure~\ref{fig:expnet}, we generate the $t$-frame expression coefficients from an audio window $a_{\{1,..,t\}}$, where the audio feature of each frame is a 0.2s mel-spectrogram. For training, we first design a ResNet-based audio encoder $\Phi_{A}$\cite{resnet, wav2lip} to embed the audio feature to a latent space. Then, a linear layer is added as the mapping network $\Phi_{M}$ to decode the expression coefficients. Here, we also add the reference expression $\beta_{0}$ from the reference image to reduce the identity uncertainty as discussed above. Since 
we use the lip-only coefficients as ground truth in the training, we explicitly add a blinking control signal $z_{blink} \in [0, 1]$ and the corresponding eye landmark loss to generate the controllable eye blinks. Formally, the network can be written as:
\begin{equation}
\beta_{\{1,...,t\}} = \Phi_{M}(\Phi_{A}(a_{\{1,...,t}\}), z_{blink}, \beta_{0})
\end{equation}


As for the loss function,
% added to map As introduce above, since our ground truth only contains the lip motions,  add a blink control signal to the mapping network and the first frame of expression as the reference. 
% these features are concatenated with the blink control signal and the expression of the original face $\beta_0$ to a linear mapping network. It contains a linear layer and aims to learn the residual expression coefficients of the input. 
we first use $\mathcal{L}_{distill}$ to evaluate the differences between the lip only expression coefficients ${R}_{e}(\mathtt{Wav2Lip}(I_0, a_{\{1,...,t\}}))$ and the generated $\beta_{\{1,...,t\}}$. Notice that, we only use the first frame $I_{0}$ of the wav2lip to generate the lip-sync video which reduces the influence of the pose variant and  other facial expressions apart from lip movement. 
Besides, we also involve the differentiable 3D face render $R_{d}$ to calculate the additional perceptual losses in explicit facial motions space.  
As shown in Figure~\ref{fig:expnet}, we calculate the landmark loss $\mathcal{L}_{lks}$ to measure the range of eye blink and the overall expression accuracy. A pretrained lip reading network $\Phi_{reader}$ is also used as temporal lip reading loss $\mathcal{L}_{read}$ to keep the perceptual lip qualities~\cite{wav2lip, spectre}. We provide more training details in the supplementary materials.

% Besides, we also consider the generated face in visual domain, several perceptual losses in the rendered face are applied to ensure the accurate expressions. \ie, the lip sync loss~\cite{wav2lip} and lip reading loss. We find that both initial face expression and the blinking signal are important to learn the accurate expression coefficient. On the one hand, the initial expression gives a truth-worth conditions which release the difficulties of this task. On the other hand, the blinking is rarely shown in the dataset which is easy to be smoothed by the network. 
\begin{figure}[tp]
    \centering
    \includegraphics[width=\columnwidth]{figs/images/pose_net.pdf}
    \vspace{-2em}
    \caption{The pipeline of the proposed PoseVAE. We learn the residual of the input head pose $\rho_{0}$ via a conditional VAE structure. Given the conditions: first frame $\rho_{0}$, style identity $Z_{style}$ and the audio clip $a_{\{1,...,t\}}$, our method learns a distribution of the residual head pose $\Delta\rho_{\{1,...,t\}} = \rho_{\{1,...,t\}} -\rho_{0}$. After training, we can generate the stylized results through the pose decoder and the conditions~$(cond.)$ only. 
    %Since we learn the relative motion of the first frame, the longer video can be generated through the head pose of the previous generated clip as shown in test part.
    }
    \label{fig:posenet}
\end{figure}

\paragraph{PoseVAE}
As shown in Figure~\ref{fig:posenet}, a VAE~\cite{vae} based model is designed to learn the realistic and identity-aware stylized head movement $ \rho \in \mathbb{R}^6$ of the real talking video. In training, the pose VAE is trained on fixed $n$ frames using an encoder-decoder-based structure. Both the encoder and decoder are two-layer MLPs, where the inputs contain a sequential $t$-frame head poses and we embed it to a Gaussian distribution. In the decoder, the network is learned to generate the $t$-frame poses from the sampled distribution. Instead of generating the pose directly, our PoseVAE learns the \textit{residual} of the condition pose $\rho_0$ of the first frame, which enables our method to generate longer, stable, and continuous head motion in testing under the condition of the first frame. 
%As shown in the testing part of Figure~\ref{fig:posenet}, the head pose of $\rho_{\{t+1, ..., 2t\}}$ can be generated by the conditional pose $\rho_{t}$.
Besides, according to CVAE~\cite{cvae}, we add the corresponding audio feature $a_{\{1, ..., t\}}$ and style identity $Z_{style}$ as conditions for rhythm awareness and identity style. The KL-divergence $\mathcal{L}_{KL}$ is used to measure the distribution of the generated motions. The mean square loss $\mathcal{L}_{MSE}$ and adversarial loss $\mathcal{L}_{GAN}$ are used to ensure the generated quality. We provide more details about the loss function in the supplementary materials.
% Finally, the total loss function can be written as:
% \begin{equation}
%     \alpha \mathcal{L}_{KL} + \sum_{i \in [1, N]}[\beta \mathcal{L}_{MSE}(\Delta \rho'_{i}, \Delta \rho_{i}) + \gamma \mathcal{L}_{GAN}(\Delta \rho'_{i}, \Delta \rho_{i})],
% \end{equation}
% where $\alpha, \beta, \gamma$ are hyper-parameters and we experimentally set them as \xiaodong{add number}, respectively.

% As defined in 3DMM, head pose $ p \in \mathbb{R}^{ 6 \times 1 } $ is a rigid transformation matrix which represents the global translation and rotation of the 3D face model. \xiaodong{todo}



\subsection{3D-aware Face Render}
\label{sec:render}
After generating the realistic 3D motion coefficients, we render the final video through a well-designed 3D-aware image animator. We get inspiration from the recent image animation method face-vid2vid~\cite{facevid2vid} because it implicitly learns the 3D information from a single image. 
However, a real video is required as the motions driving signal in their method. Our face render makes it 
 drivable through 3DMM coefficients.
As shown in Figure~\ref{fig:render}, we propose mappingNet to learn the relationship between the explicit 3DMM motion coefficients~(head pose and expression) and the implicit unsupervised 3D keypoints. Our mappingNet is built via several 1D convolutional layers. We use the temporal coefficients from a time window for smoothing as PIRenderer\cite{ren2021pirenderer}. Differently, we find the face alignment motion coefficients in PIRenderer will hugely influence the motion naturalness of audio-driven video generation and provide an experiment in Sec.~\ref{ablation}. We only use the coefficients of expression and head pose.

\begin{figure}[tp]
    \centering
    \includegraphics[width=\columnwidth]{figs/images/render.pdf}
    \vspace{-2em}
    \caption{The proposed FaceRender and comparison with face-vid2vid~\cite{facevid2vid}. Given source image $I_{s}$ and driving image $I_{d}$, face-vid2vid generates the motions in a unsupervised 3D keypoint spaces of $\mathcal{X}_{c}$, $\mathcal{X}_{s}$ and $\mathcal{X}_{d}$. Then, the image can be generated via the appearance $A_{0}$ and the keypoints. Since we do not have driving image, we use the explicit disentangled 3DMM coefficients as proxy and map it to the unsupervised 3D keypoints space.}
    \label{fig:render}
\end{figure}

As for training, our method contains two steps. Firstly, we train face-vid2vid~\cite{facevid2vid} in a self-supervised fashion as in the original paper. In the second step, we freeze all the parameters of the appearance encoder, canonical keypoints estimator, and image generator for tuning. Then, we train the mapping net on the 3DMM coefficients of the ground truth video in a reconstruction style. We give the supervision in the domain of unsupervised keypoints using $\mathcal{L}_1$ loss and the final generated video following their original implementation. More details can be founded in the supplementary materials.
% We get inspirations from recent video-driven face reenactment methods of PIRenderer~\cite{ren2021pirenderer}. PIRenderer is also a 3DMM-based face render network, which is trained to produce the dense flow fields by implicitly modulation of ADAIN~\cite{adain}. 
% However, dense flow makes it hard to preserve the identity information and it may produce obverse artifacts when the head pose is large. To this end, we   Face-vid2vid~\cite{facevid2vid} is a 3d-aware face reenactment network which learns the dense flow from a unsupversied keypoints domain but they need to be driven by a real video. Different from previous two method, we learn to map the 3DMM coefficients to the unsupervised keypoints domain of the face-vid2vid, getting a semantic disentangled face render to support our framework.

% \subsection{Loss function and training details} 
% We notice that there is a transform coefficient in the original PIRenderer, which makes the network learn the aligned face. we remove it in out coefficients.
